\subsubsection{Every hypothesis on one page}

\Cref{fig:s1forest} puts all eight preregistered comparisons side by side, cost on the left and quality on the right, so the reader can see at once which results move money and which are guards.

\begin{figure}[htbp]
\centering
\pgfplotslegendfromname{sforestlegend}\par\smallskip
% S1: every preregistered hypothesis, cost panel (log) and quality panel.
\begin{tikzpicture}
\begin{axis}[benchmark, name=cost, xmode=log, log ticks with fixed point, xtick={0.6,0.8,1,1.2}, xmin=0.5, xmax=1.4,
  scale only axis, width=5.0cm, height=6.2cm, xmajorgrids, ymin=-0.6, ymax=7.95, ytick={0,...,7},
  yticklabels from table={generated/data/s1-forest.dat}{label}, y tick label style={font=\footnotesize},
  xlabel={Cost ratio (log)}, title={Cost}, title style={font=\small\bfseries},
  legend to name=sforestlegend, legend columns=3, legend cell align=left,
  legend style={/tikz/every even column/.append style={column sep=10pt}}]
\draw[black!70, line width=0.8pt] (axis cs:1,-0.6) -- (axis cs:1,7.6);
\draw[okorange, dashed, line width=0.8pt] (axis cs:\SThreshOne,-0.6) -- (axis cs:\SThreshOne,7.6);
\addplot[nodes near coords, point meta=explicit symbolic, nodes near coords style={font=\scriptsize, fill=white, fill opacity=0.85, text opacity=1, inner sep=1pt, anchor=south, yshift=2.5pt, text=black!80}, only marks, mark=*, mark size=2.4pt, color=okblue, error bars/.cd, x dir=both, x explicit,
  error bar style={line width=1pt}] table[meta=v, x=r, y=y, x error minus=rm, x error plus=rp]{generated/data/s1-forest.dat};
\addlegendentry{estimate, 95\,\% CI}
\addlegendimage{black!70, line width=0.8pt}\addlegendentry{no difference}
\addlegendimage{okorange, dashed, line width=0.8pt}\addlegendentry{preregistered threshold}
\end{axis}
\begin{axis}[benchmark, at={(cost.south east)}, xshift=0.5cm, scale only axis, width=4.0cm, height=6.2cm,
  xmin=-0.09, xmax=0.06, xtick={-0.05,0,0.05}, scaled x ticks=false, xticklabel style={/pgf/number format/fixed}, xmajorgrids,
  ymin=-0.6, ymax=7.95, ytick={0,...,7}, yticklabels={}, xlabel={Turn-pass difference}, title={Quality},
  title style={font=\small\bfseries}]
\draw[black!70, line width=0.8pt] (axis cs:0,-0.6) -- (axis cs:0,7.6);
\draw[okorange, dashed, line width=0.8pt] (axis cs:-0.05,-0.6) -- (axis cs:-0.05,7.6);
\addplot[only marks, mark=square*, mark size=2.2pt, color=okgreen, error bars/.cd, x dir=both, x explicit,
  error bar style={line width=1pt}] table[x=d, y=y, x error minus=dm, x error plus=dp]{generated/data/s1-forest.dat};
\end{axis}
\end{tikzpicture}

\caption{Two hypotheses carry money (H1, H5); the rest are guards. S1 hypotheses H1--H7: cost ratio (left, log scale) and turn-pass difference (right), each with its 95\,\%
interval. H3b is descriptive. Thresholds: the H1 cost bound (\SThreshOne) and the quality margin (\NIMargin).}
\label{fig:s1forest}
\howtoread{Read each row across both panels. A row helps a configuration when its cost dot is left of 1 and its
quality whisker stays right of the dashed line. H1 does both. H5 is far right of 1: routing on Opus costs more. H4
straddles 1: medium effort did nothing measurable on Opus. H7's quality whisker crosses the dashed line, which is why
review and explain stay on the host.}
\end{figure}

\begin{keyidea}
Two of the eight rows carry money: H1 (Fable routing saves about two fifths) and H5 (Opus routing loses about a
quarter). The rest are guards: the bundle is cost-equivalent within the \(\pm\)5\,\% margin (H3), the model adds nothing over the rule (H2), live sessions behave as
predicted (H6), and review and explain need protection (H7).
\end{keyidea}

\subsubsection{How to read a paired cost ratio}\label{sec:worked}

\begin{casecard}{Worked example: one scenario, two reps, one ratio}
Take one explain scenario from the holdout study, the scenario whose ratio sits closest to the overall result. Under the frozen Fable
configuration it was \WxArm. Across its \WxReps\ repetitions the routed sessions cost \$\WxRouted\ on average and the
plain Fable sessions \$\WxPlain. Its paired cost ratio is therefore $\WxRouted / \WxPlain = \WxRatio$; on the log scale
that is $\ln \WxRatio$ = \WxLog. Doing the same for each of the \WxNScen\ scenarios, averaging the logs and
exponentiating gives the \term{geometric mean}, \WxGM: the H1 headline. The interval comes from resampling whole
scenarios \SResamples\ times and repeating the calculation.
\end{casecard}

Why a geometric mean? A ratio of 0.5 and a ratio of 2 should cancel; their arithmetic mean (1.25) says ``more
expensive'', their geometric mean (1) says ``no change''. Averaging logs treats halving and doubling symmetrically.

\subsubsection{The variance underneath the headline}

A geometric mean hides how much individual tasks vary. \Cref{fig:perscen} shows every scenario's paired ratio, sorted, for the two results that carry money.

\begin{figure}[htbp]
\centering
\pgfplotslegendfromname{pslegend}\par\smallskip
% S1: per-scenario paired cost ratios, sorted (log y).
\begin{tikzpicture}
\begin{axis}[benchmark, name=f, ymode=log, log ticks with fixed point, ytick={0.4,0.6,0.8,1,1.5,2}, ymin=0.33, ymax=2.4,
  scale only axis, width=5.6cm, height=5cm, ymajorgrids, xmin=0, xmax=61, xlabel={Scenario (sorted)},
  ylabel={Cost ratio (log)}, title={Fable: frozen config / plain}, title style={font=\small\bfseries},
  legend to name=pslegend, legend columns=2, legend cell align=left,
  legend style={/tikz/every even column/.append style={column sep=10pt}}]
\draw[black!70, line width=0.8pt] (axis cs:0,1) -- (axis cs:61,1);
\addplot[only marks, mark=*, mark size=1.6pt, color=okblue] table[x=x, y=r]{generated/data/perscen-fable.dat};
\addlegendentry{one scenario (mean of its reps)}
\addplot[only marks, mark=triangle*, mark size=1.9pt, color=okorange] coordinates {(-10,1)};
\addlegendentry{Opus panel}
\end{axis}
\begin{axis}[benchmark, at={(f.south east)}, xshift=1.2cm, ymode=log, log ticks with fixed point,
  ytick={0.4,0.6,0.8,1,1.5,2}, ymin=0.33, ymax=2.4, scale only axis, width=5.6cm, height=5cm, ymajorgrids,
  xmin=0, xmax=61, xlabel={Scenario (sorted)}, title={Opus: Sonnet / plain}, title style={font=\small\bfseries}]
\draw[black!70, line width=0.8pt] (axis cs:0,1) -- (axis cs:61,1);
\addplot[only marks, mark=triangle*, mark size=1.9pt, color=okorange] table[x=x, y=r]{generated/data/perscen-opus.dat};
\end{axis}
\end{tikzpicture}

\caption{The direction is consistent across scenarios; the size varies widely. Per-scenario paired cost ratios in S1, sorted. Left: the frozen Fable configuration against plain Fable
(\PsFableN\ scenarios). Right: Sonnet decided once against plain Opus (\PsOpusN\ scenarios). Log scale.}
\label{fig:perscen}
\howtoread{Each dot is one scenario. Below the line saves money. On Fable, \PsFableAbove\ of \PsFableN\ scenarios cost
more than plain Fable, and the ratios run from \PsFableMin\ to \PsFableMax. On Opus, \PsOpusAbove\ of \PsOpusN\ cost
more, from \PsOpusMin\ to \PsOpusMax. The direction is consistent; the size varies a lot from task to task.}
\end{figure}

\subsubsection{Where the risk sits: task types and families}

The quality cost of routing is not spread evenly. \Cref{fig:subgroups} splits the Fable result by task type and by scenario family (exploratory).

\begin{figure}[htbp]
\centering
\pgfplotslegendfromname{sublegend}\par\smallskip
% S1, Fable: the frozen configuration against plain Fable, per task type and per family.
\begin{tikzpicture}
\begin{axis}[benchmark, name=cost, xmode=log, log ticks with fixed point, xtick={0.4,0.6,0.8,1}, xmin=0.35, xmax=1.1,
  scale only axis, width=5.0cm, height=6.6cm, xmajorgrids, ymin=-0.6, ymax=10.95, ytick={\SubTicks},
  yticklabels/.expanded={\SubLabels}, y tick label style={font=\footnotesize},
  xlabel={Cost ratio (log)}, title={Cost}, title style={font=\small\bfseries},
  legend to name=sublegend, legend columns=2, legend cell align=left,
  legend style={/tikz/every even column/.append style={column sep=10pt}}]
\draw[black!70, line width=0.8pt] (axis cs:1,-0.6) -- (axis cs:1,10.6);
\addplot[nodes near coords, point meta=explicit symbolic, nodes near coords style={font=\scriptsize, fill=white, fill opacity=0.85, text opacity=1, inner sep=1pt, anchor=south, yshift=2.5pt, text=black!80}, only marks, mark=*, mark size=2.4pt, color=okblue, error bars/.cd, x dir=both, x explicit,
  error bar style={line width=1pt}] table[meta=v, x=r, y=y, x error minus=rm, x error plus=rp]{generated/data/s1-subgroups.dat};
\addlegendentry{estimate, 95\,\% CI}
\addlegendimage{okorange, dashed, line width=0.8pt}\addlegendentry{quality margin}
\end{axis}
\begin{axis}[benchmark, at={(cost.south east)}, xshift=0.5cm, scale only axis, width=4.0cm, height=6.6cm,
  xmin=-0.12, xmax=0.07, xtick={-0.1,-0.05,0,0.05}, scaled x ticks=false, xticklabel style={/pgf/number format/fixed}, xmajorgrids,
  ymin=-0.6, ymax=10.95, ytick={\SubTicks}, yticklabels={}, xlabel={Turn-pass difference}, title={Quality},
  title style={font=\small\bfseries}]
\draw[black!70, line width=0.8pt] (axis cs:0,-0.6) -- (axis cs:0,10.6);
\draw[okorange, dashed, line width=0.8pt] (axis cs:-0.05,-0.6) -- (axis cs:-0.05,10.6);
\addplot[only marks, mark=square*, mark size=2.2pt, color=okgreen, error bars/.cd, x dir=both, x explicit,
  error bar style={line width=1pt}] table[x=d, y=y, x error minus=dm, x error plus=dp]{generated/data/s1-subgroups.dat};
\end{axis}
\end{tikzpicture}

\caption{Every subgroup saves money; quality intervals for several subgroups are too wide to rule out a loss. Lower bounds fall below the $-$0.05 margin for the explain (\SubExplainTaskLow), mixed (\SubMixedTaskLow) and docs (\SubDocsTaskLow) tasks and the mixed (\SubMixedFamilyLow) and knowledge (\SubKnowledgeFamilyLow) families; review's lower bound is \SubReviewTaskLow. Subgroups hold 9--20 scenarios; treat as hypothesis-generating. The frozen Fable configuration against plain Fable, split by task type and by scenario family (number of
cost-valid scenarios in brackets; the quality panel is unfiltered and uses every scenario of the subgroup), with 95\,\%
intervals. Exploratory: subgroups were not preregistered except H7.}
\label{fig:subgroups}
\howtoread{Every subgroup saves money. The quality panel is where they differ: explain (\SubExplainTaskDtp; lower
bound \SubExplainTaskLow) and mixed work cross the margin, review touches it (\SubReviewTaskLow), while bugfix and
feature sit near zero. Small subgroups give wide whiskers; treat single rows with caution.}
\end{figure}

\subsubsection{Which configuration the rule picked}

The freeze rule considered every configuration the preregistration listed and chose one per host. \Cref{fig:fz} shows them all, with the choice and the preregistered configuration marked; Table~\ref{tab:fzkey} lists them.

\begin{figure}[htbp]
\centering
\pgfplotslegendfromname{fzlegendFable}\par\smallskip
\fzpanel{Fable}{Fable}{fable}\par\medskip
\fzpanel{Opus}{Opus}{opus}
\caption{On Fable many configurations qualify; on Opus none does. All \FzFableN\ Fable and \FzOpusN\ Opus configurations the freeze rule considered: cost ratio against the
host's plain default and turn-pass difference (point estimates). Keys group configurations that land on the same
point; Table~\ref{tab:fzkey} lists them.}
\label{fig:fz}
\howtoread{Left and up is better. Filled dots passed the preregistered candidate test (cost upper bound below 1,
quality lower bound above the margin, final-state point no worse than the specified limit). The star is the rule's choice; the
diamond is the preregistered configuration. On Fable \FzFableCand\ configurations qualified; the cheapest cluster
drops the scope gate. On Opus none qualified, so the shipped default stays.}
\end{figure}

\subsubsection{The decision model against the rule}\label{sec:bound}

Jev and the rule R* chose differently on only \DcN\ of \DcUnits\ Fable scenario-reps (\cref{tab:discordant}).

\begin{table}[htbp]
\centering\small
\caption{The \DcN\ S1 scenario-reps on which the decision model (Jev) and the rule R* chose different session
models, with what each choice cost and how it scored.}
\label{tab:discordant}
\begin{tabular*}{\textwidth}{@{\extracolsep{\fill}}l r l l r r r r@{}}
\toprule
 & & \multicolumn{2}{c}{model chosen by} & \multicolumn{2}{c}{session cost} & \multicolumn{2}{c}{turn-pass} \\
\cmidrule(lr){3-4}\cmidrule(lr){5-6}\cmidrule(l){7-8}
Scenario & rep & Jev & rule R* & Jev & R* & Jev & R* \\
\midrule
go-bowling & 1 & Fable, medium & Sonnet (Pc) & \$6.53 & \$4.06 & 0.857 & 1.000 \\
go-bowling & 2 & Fable, medium & Sonnet (Pc) & \$5.66 & \$4.77 & 0.929 & 0.929 \\
peewee-filterfk & 1 & Fable, medium & Sonnet (Pc) & \$2.15 & \$1.94 & 0.889 & 1.000 \\
peewee-filterfk & 2 & Fable, medium & Sonnet (Pc) & \$2.34 & \$1.87 & 0.889 & 1.000 \\
\bottomrule
\end{tabular*}

\end{table}

\begin{casecard}{Worked example: why the two deciders cannot differ much}
Two deciders that agree on a scenario-rep produce the same session there, so their difference comes only from the
scenario-reps where they disagree. The share of those is $d = \DcN/\DcUnits = \DcD$. On them the choices differed in
cost by \$\DcMeanAbsDiff\ per session on average. So the most the deciders can differ, averaged over all sessions, is
$d \times \$\DcMeanAbsDiff = \$\DcBound$ per session, about \DcBoundPct\,\% of a typical session (\$\DcMeanCost).
That bound, not a significance test, is what makes ``equivalent'' believable: with so few disagreements, no amount of
skill on those few can move the total.
\end{casecard}

\subsubsection{Effort, across three hosts}

Three separate studies measured medium against default reasoning effort, one per host model. \Cref{fig:effort} puts them on one chart.

\begin{figure}[htbp]
\centering
\pgfplotslegendfromname{efflegend}\par\smallskip
% Medium vs default reasoning effort, per host, three studies.
\begin{tikzpicture}
\begin{axis}[benchmark, name=cost, scale only axis, width=5.2cm, height=2.8cm, xmin=0.75, xmax=1.05,
  xtick={0.8,0.9,1}, xmajorgrids, ymin=-0.6, ymax=2.95, ytick={0,1,2},
  yticklabels from table={generated/data/effort-hosts.dat}{label}, y tick label style={font=\footnotesize},
  xlabel={Cost ratio, medium / default}, title={Cost}, title style={font=\small\bfseries},
  legend to name=efflegend, legend columns=2, legend cell align=left,
  legend style={/tikz/every even column/.append style={column sep=10pt}}]
\draw[black!70, line width=0.8pt] (axis cs:1,-0.6) -- (axis cs:1,2.6);
\addplot[nodes near coords, point meta=explicit symbolic, nodes near coords style={font=\scriptsize, fill=white, fill opacity=0.85, text opacity=1, inner sep=1pt, anchor=south, yshift=2.5pt, text=black!80}, only marks, mark=*, mark size=2.4pt, color=okblue, error bars/.cd, x dir=both, x explicit,
  error bar style={line width=1pt}] table[meta=v, x=r, y=y, x error minus=rm, x error plus=rp]{generated/data/effort-hosts.dat};
\addlegendentry{cost, 95\,\% CI}
\addplot[only marks, mark=square*, mark size=2.2pt, color=okgreen] coordinates {(5,5)};
\addlegendentry{turn-pass, 95\,\% CI}
\end{axis}
\begin{axis}[benchmark, at={(cost.south east)}, xshift=0.5cm, scale only axis, width=4.0cm, height=2.8cm,
  xmin=-0.06, xmax=0.09, xtick={-0.05,0,0.05}, scaled x ticks=false, xticklabel style={/pgf/number format/fixed}, xmajorgrids,
  ymin=-0.6, ymax=2.95, ytick={0,1,2}, yticklabels={}, xlabel={Turn-pass difference}, title={Quality},
  title style={font=\small\bfseries}]
\draw[black!70, line width=0.8pt] (axis cs:0,-0.6) -- (axis cs:0,2.6);
\draw[okorange, dashed, line width=0.8pt] (axis cs:-0.05,-0.6) -- (axis cs:-0.05,2.6);
\addplot[only marks, mark=square*, mark size=2.2pt, color=okgreen, error bars/.cd, x dir=both, x explicit,
  error bar style={line width=1pt}] table[x=d, y=y, x error minus=dm, x error plus=dp]{generated/data/effort-hosts.dat};
\end{axis}
\end{tikzpicture}

\caption{Medium effort saved money on Sonnet and Fable but did not demonstrate a saving on Opus. Medium against default reasoning effort, measured in three studies: Sonnet (follow-up 1), Fable (follow-up
2) and Opus (S1, H4). 95\,\% intervals.}
\label{fig:effort}
\howtoread{Sonnet and Fable sit well left of 1: medium effort is cheaper with no detectable quality loss. Opus sits
on 1. The same setting saves money on two models and nothing on the third, so effort must be set per host.}
\end{figure}

\subsubsection{Where the money goes}

Why does the same switch save money on one host and cost money on the other? \Cref{fig:composition} splits the average S1 session's cost by token class.

\begin{figure}[htbp]
\centering
\pgfplotslegendfromname{complegend}\par\smallskip
% S1: mean dollars per session by token class.
\begin{tikzpicture}
\begin{axis}[benchmark, xbar stacked, bar width=10pt, scale only axis, width=10cm, height=4.2cm, xmin=0, xmajorgrids,
  ymin=-0.6, ymax=3.6, ytick={0,1,2,3}, yticklabels from table={generated/data/s1-composition.dat}{label},
  y tick label style={font=\small}, xlabel={Mean cost per session (US dollars)},
  title={Cost per session, by token class}, title style={font=\small\bfseries},
  legend to name=complegend, legend columns=4, legend cell align=left,
  legend style={/tikz/every even column/.append style={column sep=8pt}}]
\addplot[fill=okgray!60, draw=black!60] table[x=inp, y=y]{generated/data/s1-composition.dat};
\addlegendentry{uncached input}
\addplot[fill=okblue!55, draw=black!60] table[x=rd, y=y]{generated/data/s1-composition.dat};
\addlegendentry{cache read}
\addplot[fill=okorange!80, draw=black!60, postaction={pattern=north east lines, pattern color=white}] table[x=wr, y=y]{generated/data/s1-composition.dat};
\addlegendentry{cache write}
\addplot[fill=okgreen!60, draw=black!60, postaction={pattern=dots, pattern color=white}] table[x=out, y=y]{generated/data/s1-composition.dat};
\addlegendentry{output}
\node[anchor=west, font=\scriptsize, text=black!80] at (axis cs:1.9122,0) {1.91};
\node[anchor=west, font=\scriptsize, text=black!80] at (axis cs:2.3392,1) {2.34};
\node[anchor=west, font=\scriptsize, text=black!80] at (axis cs:3.0271,2) {3.03};
\node[anchor=west, font=\scriptsize, text=black!80] at (axis cs:4.5530,3) {4.55};

\end{axis}
\end{tikzpicture}

\caption{Plain Fable spends most of its money writing to the cache. Mean cost per S1 session split by token class, at the study's price table (exploratory).}
\label{fig:composition}
\howtoread{Each bar is the average session. Plain Fable spends most on cache writes (\CcPlainFableWriteShare\,\% of
\$\CcPlainFableTotal): writing context into Fable's cache is expensive. The Sonnet session (\$\CcSonnetDecidedOTotal)
spends more on cache reads (\CcSonnetDecidedOReadShare\,\%), which are cheap per token but many. Plain Opus
(\$\CcPlainOpusTotal) is already cheaper than the Sonnet session, which is the whole H5 story in one bar.}
\end{figure}


\subsubsection{Every configuration the freeze rule considered}


\begingroup\small
\FloatBarrier
\begin{xltabular}{\textwidth}{l l >{\raggedright\arraybackslash}X S[table-format=1.3] S[table-format=-1.3] c}
\caption{Every configuration the S1 freeze rule considered (keys as in \cref{fig:fz}). Cost ratio and turn-pass difference against the host's plain default.}\label{tab:fzkey}\\
\toprule Key & Host & Configuration & {cost ratio} & {turn-pass $\Delta$} & candidate \\ \midrule \endfirsthead
\toprule Key & Host & Configuration & {cost ratio} & {turn-pass $\Delta$} & candidate \\ \midrule \endhead
F1 & Fable & always\_route | strong default | scope off | keep none & 0.528 & -0.017 & \ok \\
F1 & Fable & always\_route | strong medium | scope off | keep none & 0.528 & -0.017 & \ok \\
F1 & Fable & rstar | strong default | scope off | keep none & 0.528 & -0.017 & \ok \\
F1 & Fable & rstar | strong medium | scope off | keep none & 0.528 & -0.017 & \ok \\
F2 & Fable & always\_route | strong medium | scope 300 | keep none & 0.581 & -0.013 & \ok \\
F2 & Fable & rstar | strong medium | scope 300 | keep none & 0.581 & -0.013 & \ok \\
F2 & Fable & jev | strong medium | scope 300 | keep none & 0.586 & -0.016 & \ok \\
F3 & Fable & always\_route | strong default | scope 300 | keep none & 0.608 & -0.012 & \ok \\
F3 & Fable & rstar | strong default | scope 300 | keep none & 0.608 & -0.012 & \ok \\
F3 & Fable & jev | strong default | scope 300 | keep none & 0.618 & -0.015 & \ok \\
F4 & Fable & always\_route | strong medium | scope off | keep review+explain & 0.618 & -0.007 & \ok \\
F4 & Fable & rstar | strong medium | scope off | keep review+explain & 0.618 & -0.007 & \ok \\
F5 & Fable & always\_route | strong medium | scope 300 | keep review+explain & 0.650 & -0.006 & \ok \\
F5 & Fable & rstar | strong medium | scope 300 | keep review+explain & 0.650 & -0.006 & \ok \\
F5 & Fable & jev | strong medium | scope 300 | keep review+explain & 0.655 & -0.009 & \ok \\
F6 & Fable & always\_route | strong default | scope off | keep review+explain & 0.650 & -0.001 & \ok \\
F6 & Fable & rstar | strong default | scope off | keep review+explain & 0.650 & -0.001 & \ok \\
F7 & Fable & always\_route | strong default | scope 300 | keep review+explain & 0.706 & -0.000 & \ok \\
F7 & Fable & rstar | strong default | scope 300 | keep review+explain & 0.706 & -0.000 & \ok \\
F7 & Fable & jev | strong default | scope 300 | keep review+explain & 0.717 & -0.003 & \ok \\
F8 & Fable & plain-medium & 0.851 & 0.005 & \ok \\
F9 & Fable & bundle-unrouted medium & 0.853 & -0.005 & \ok \\
F10 & Fable & bundle-unrouted default & 1.004 & 0.004 &  \\
O1 & Opus & plain-medium & 0.992 & -0.002 &  \\
O2 & Opus & bundle-unrouted default & 0.999 & 0.007 &  \\
O2 & Opus & bundle-unrouted medium & 1.009 & 0.003 &  \\
O3 & Opus & jev | strong default | scope 300 | keep review+explain & 1.080 & 0.005 &  \\
O4 & Opus & jev | strong medium | scope 300 | keep review+explain & 1.091 & 0.000 &  \\
O5 & Opus & always\_route | strong default | scope 300 | keep review+explain & 1.097 & 0.005 &  \\
O5 & Opus & rstar | strong default | scope 300 | keep review+explain & 1.097 & 0.005 &  \\
O6 & Opus & always\_route | strong medium | scope 300 | keep review+explain & 1.109 & 0.004 &  \\
O6 & Opus & rstar | strong medium | scope 300 | keep review+explain & 1.109 & 0.004 &  \\
O7 & Opus & jev | strong default | scope 300 | keep none & 1.136 & -0.002 &  \\
O7 & Opus & jev | strong medium | scope 300 | keep none & 1.144 & -0.006 &  \\
O8 & Opus & always\_route | strong default | scope 300 | keep none & 1.153 & -0.002 &  \\
O8 & Opus & rstar | strong default | scope 300 | keep none & 1.153 & -0.002 &  \\
O8 & Opus & always\_route | strong medium | scope 300 | keep none & 1.163 & -0.002 &  \\
O8 & Opus & rstar | strong medium | scope 300 | keep none & 1.163 & -0.002 &  \\
O9 & Opus & always\_route | strong default | scope off | keep review+explain & 1.182 & 0.005 &  \\
O9 & Opus & rstar | strong default | scope off | keep review+explain & 1.182 & 0.005 &  \\
O9 & Opus & always\_route | strong medium | scope off | keep review+explain & 1.187 & 0.005 &  \\
O9 & Opus & rstar | strong medium | scope off | keep review+explain & 1.187 & 0.005 &  \\
O10 & Opus & always\_route | strong default | scope off | keep none & 1.255 & -0.004 &  \\
O10 & Opus & always\_route | strong medium | scope off | keep none & 1.255 & -0.004 &  \\
O10 & Opus & rstar | strong default | scope off | keep none & 1.255 & -0.004 &  \\
O10 & Opus & rstar | strong medium | scope off | keep none & 1.255 & -0.004 &  \\
\bottomrule
\end{xltabular}

\FloatBarrier
\endgroup

